\section{Parte II Reporte del proyecto}
\subsection{Panorama y arquitectura}
El caso de uso es un agente que obtiene datos, los limpia, calcula
estadísticas, detecta anomalías, produce un histograma y compone un
resumen. Se busca menor latencia conservando todos los insumos y la
respuesta de referencia. El límite inicial es $p$; se observa CPU
y RSS individual, sin imponer aún cuotas de memoria.
\begin{center}\small
Plan JSON $\rightarrow$ análisis $\rightarrow$ DAG ponderado\\
$\rightarrow$ rank 0 $\rightleftarrows$ workers MPI\\
$\rightarrow$ herramientas MCP $\rightarrow$ resultados e historial.
\end{center}
Rank 0 resuelve referencias y recibe con \texttt{ANY\_SOURCE}.
Cada worker ejecuta un proceso aislado. El cliente MCP inicializa
la sesión stdio, invoca la herramienta, valida su estado y retorna
resultados estructurados. Los workers no necesitan el grafo completo.

\subsection{M1 Algoritmo central implementado}
Se implementaron ambas políticas, conflictos, reintentos, timeout,
exportación DOT y reportes JSON/HTML. SQLite aprende
$\hat w_{\rm nueva}=0.3w_{\rm observado}+0.7\hat w_{\rm anterior}$.
Los pesos explícitos prevalecen sobre historial, intervalo Jev y
heurística. Solo se aprende de éxitos y no se cambian prioridades
a mitad del ensayo.

El plan original de siete tareas completó MPI con rank 0 y cuatro
workers. Cuatro herramientas MCP se ejecutaron con rank 0 y tres workers,
conservando el resultado local. Se probaron Python 3.12.4, mpi4py 4.1.2,
MPICH 5.0.2 y MCP SDK 1.30.0. Se verificó stdio y varios procesos
en una computadora; no se afirma validación multinodo.

\begin{table}[!t]
\caption{Evaluación local. Medianas de tres repeticiones, en segundos.
FIFO y crítica usan dos workers; secuencial, uno.}
\label{tab:local}
\centering\small
\begin{tabular}{lrrrr}
\toprule
Carga & Tareas & Sec. & FIFO & Crítica\\
\midrule
Consulta uniforme & 5 & \UniformSeq & \UniformFifo & \UniformCrit\\
Consulta constante & 5 & \ConstantSeq & \ConstantFifo & \ConstantCrit\\
Consulta con extremo & 5 & \OutlierSeq & \OutlierFifo & \OutlierCrit\\
Rama crítica & 4 & \PrioritySeq & \PriorityFifo & \PriorityCrit\\
Abanico pequeño & 9 & \ScaleSmallSeq & \ScaleSmallFifo & \ScaleSmallCrit\\
Abanico mayor & 25 & \ScaleLargeSeq & \ScaleLargeFifo & \ScaleLargeCrit\\
\bottomrule
\end{tabular}
\end{table}

La carga con rama crítica ahorró \PrioritySaving{}\% de tiempo frente
al secuencial y \PriorityFifoSaving{}\% frente a FIFO.
Las consultas pequeñas no muestran dominio consistente sobre FIFO
(Tabla~\ref{tab:local}). La carga de 25 tareas alcanzó speedup
\ScaleSpeedup{} y eficiencia \ScaleEfficiency{}\% con dos workers.
Su CPU acumulado de tareas fue \ScaleCpuSeq{} s secuencial frente a
\ScaleCpuCrit{} s con prioridad crítica: reducir latencia no demuestra
ahorrar trabajo computacional ni energía.

\subsection{M2 Robustez y evaluación determinista implementadas}
Se fijaron tres consultas en \texttt{experiments/golden.json}:
$[1,2,3,4,5]$, $[7,7,7,7]$ y $[10,20,30,40,100]$.
Sus referencias se calcularon directamente: estadística poblacional,
anomalías con $|z|>1.5$ y cinco bins incluyendo el extremo superior.
La tercera tiene media 40, mediana 30, desviación $\sqrt{1000}$,
extremo en índice 4 y frecuencias $[2,2,0,0,1]$.

Los 27 ensayos de consultas alcanzaron exactitud, completitud y corrección
de precedencia 1.00, con hashes iguales entre métodos y repeticiones.
La carga de prioridad y dos abanicos completan 54 ejecuciones, manteniendo
valores, herramientas, pesos e intérprete, sin historial. Se registraron
mediana, mínimo, máximo, CPU, RSS y hashes; los tiempos no son deterministas.

El proveedor fue heurístico, sin LLM ni muestreo: tokens remotos
y costo de API fueron cero; energía y costo del hardware no se midieron.
El mayor RSS individual fue \MaxRss{} MiB; excluye memoria agregada
y servidores MCP hijos. Hay \TestCount{} pruebas automatizadas de
contrato, análisis, ejecución, FIFO y MCP, incluyendo cancelación.
El sistema operativo puede cambiar recepción y asignación física:
se exige igualdad de resultados para herramientas puras, no cronogramas
idénticos.

Jev se probó con respuestas controladas, sin consulta autenticada.
La evaluación pendiente fijará un ID inmutable de modelo y archivará
solicitudes y respuestas, repitiendo diez análisis por plan contra
aristas etiquetadas. No se inventarán parámetros de temperatura o semilla
que la API no expone. La variabilidad del modelo se separará de la
del scheduler y de las herramientas.

\subsection{M3 Calidad y asignación de recursos planeadas}
La próxima versión minimizará $C_{\max}$ y, secundariamente, CPU y
costo remoto, con exactitud y completitud 1.00 en consultas deterministas,
cero violaciones de precedencia y acuerdo de hashes 1.00.
Para inferencia se propone precisión y recall de aristas al menos 0.99
en planes revisados. Son umbrales de aceptación, no resultados medidos.
Una relación incierta detendrá el plan o exigirá una dependencia
explícita validada.

Se seleccionará el menor $p$ que cumpla latencia y presupuesto de
memoria. Los ensayos MPI iniciales muestran por qué: las medianas
de tres repeticiones fueron \MpiTwo{} s con dos workers y
\MpiFour{} s con cuatro, frente a \MpiSeq{} s secuenciales.
Hubo dispersión y carga de fondo no aislada. Son mediciones exploratorias,
no evidencia de escalamiento monotónico; el benchmark final aislará
estos factores.

\subsection{M4 Benchmark y sistema final planeados}
Se conservará coordinador--workers, con admisión por memoria,
artefactos por referencia y sesiones MCP persistentes. Se comparará
secuencial, FIFO, prioridad crítica y una variante de HEFT con
comunicación y heterogeneidad. Kahn, listas y camino crítico son
las bases clásicas elegidas; su correspondencia exacta con el temario
del curso debe confirmarse.

La matriz prevista contiene 8, 32, 128 y 512 tareas, distinta anchura
y densidad, varios tamaños de entrada y $p\in\{1,2,4,8\}$.
Habrá calentamiento, diez o más repeticiones, orden aleatorizado,
máquina sin carga ajena y versiones congeladas. Se reportarán mediana,
percentil 95 e intervalo de confianza del speedup, verificando referencias
y hashes por ensayo. Se distinguirán pruebas multinodo de las locales.

\begin{table}[!t]
\caption{Métricas del benchmark final.}
\centering\small
\begin{tabular}{p{.23\columnwidth}p{.66\columnwidth}}
\toprule
Dimensión & Medición\\
\midrule
Eficiencia & Tiempo, CPU, tokens, RSS agregado, bytes y costo con tarifa fechada\\
Calidad & Exactitud, corrección del DAG y completitud\\
Consistencia & Acuerdo de hashes y decisiones repetidas\\
Determinismo & Configuración fija; separar resultados, orden y variación Jev\\
Escalabilidad & Tiempo y eficiencia por carga y workers\\
\bottomrule
\end{tabular}
\end{table}

Se aceptará ahorro solo si se mantienen los umbrales y mejora la
métrica objetivo. El costo de API sumará tokens por tarifa real,
sin confundirlo con energía ni CPU. La memoria total exigirá muestrear
ranks, procesos de tarea y servidores, complementando el RSS individual.

\subsection{Repositorio y hoja de ruta}
Se entregan código, planes, esquema, servidor MCP, pruebas, referencias
fijas, evaluación, resultados y fuentes del informe. Las versiones están
congeladas. El contenedor no se presenta como verificado.

M1 está implementado y verificado en MPI. M2 tiene evidencia determinista
y robustez; la variabilidad Jev está pendiente. M3 requiere costos
condicionados, cuotas y selección de $p$. M4 requiere comparadores
heterogéneos, aislamiento y validación multinodo. El orden será cerrar
aristas de referencia y fijar versión Jev, implementar admisión y
memoria total, comparar cargas ampliadas y reproducir en varios nodos.
